% Part: first-order-logic
% Chapter: axiomatic-deduction
% Section: axioms-rules-quantifiers

\documentclass[../../../include/open-logic-section]{subfiles}

\begin{document}

\olfileid{fol}{axd}{qua}

\olsection{పరిమాణీకరణికాల స్వీకృతాలు, నియమాలు}

\begin{defn}[పరిమాణీకరణికాల స్వీకృతాలు]
పరిమాణీకరణికాలకు సంబంధించిన \emph{స్వీకృతాలు} కింది రూపాల
దృష్టాంతాలన్నీ:
\begin{align}
\ollabel{ax:q1} & \lforall[x][!B] \lif !B(t), \\
\ollabel{ax:q2} & !B(t) \lif \lexists[x][!B].
\end{align}
ఇక్కడ $t$ ఏ సంవృత పదమైనా కావచ్చు.
\end{defn}

\begin{defn}[పరిమాణీకరణికాల నియమాలు]\begin{enumerate}
\item ఒక \tetoken{వ్యుత్పత్తిలో}{!!{derivation}} $!B \lif !A(a)$ అంతకు ముందే వచ్చి ఉండి, $a$
  అనేది~$\Gamma$లో గానీ~$!B$లో గానీ $!A(x)$లో గానీ కనిపించకపోతే, $!B \lif
  \lforall[x][!A(x)]$ ఒక సరైన నిగమన దశ.
\item ఒక \tetoken{వ్యుత్పత్తిలో}{!!{derivation}} $!A(a) \lif !B$ అంతకు ముందే వచ్చి ఉండి, $a$
  అనేది~$\Gamma$లో గానీ~$!B$లో గానీ $!A(x)$లో గానీ కనిపించకపోతే,
  $\lexists[x][!A(x)] \lif !B$ ఒక సరైన నిగమన దశ.
\end{enumerate}\end{defn}

ఈ రెండింటిలో దేనినైనా ``\QR'' అని సంక్షిప్తంగా రాస్తాం.

\end{document}
